\documentclass[10pt,conference]{IEEEtran}
\usepackage{graphicx}
\usepackage{booktabs}
\usepackage{amsmath}
\usepackage{url}
\title{BodyShield: Falsifying and Repairing Hidden Embodiment-Control Assumptions in Robot Policies}
\author{Anonymous Authors}
\begin{document}
\maketitle
\begin{abstract}
Robot policies can succeed in nominal evaluations while relying on brittle hidden assumptions about latency, calibration, joint range, gripper authority, sensing geometry, payload, and contact. BodyShield exposes and repairs hidden embodiment-control assumptions. This draft contains analytic-simulation evidence, generated frames rather than real camera videos, synthetic visual and trajectory proxy audits, bounded MuJoCo/ManiSkill probes, and release/package verification; external full-scale trained-policy high-fidelity benchmarks remain future evidence tiers, and real robot results are blocked until safety gates pass.
\end{abstract}
\section{Introduction}
Nominal success is a weak certificate for robot policies because the policy may depend on a hidden body/control interface detail. BodyShield is a falsification-to-repair method: BodyBreak searches for low-cost embodiment-control perturbations that break a policy, then BodyShield repairs the action representation or policy parameters against discovered failures and evaluates held-out perturbations. The goal is adaptive embodied-intelligence evidence about failure diagnosis and repair, not a cheap-arm demo or a broad cross-embodiment-transfer claim.
\section{Related Work}
Domain randomization and dynamics randomization sample broad distributions for sim-to-real transfer \cite{tobin2017domainrandomization,peng2018dynamicsrandomization,openai2018dexterous,openai2019rubiks}. Embodiment-aware deployment methods such as UMI-on-Air and EmbodiSteer steer policies toward target-body feasibility \cite{gupta2025umionair,wang2026embodisteer}. Counterexample-guided and verification-guided safe-RL methods search for violations and sometimes repair or shield policies \cite{karunakaran2020counterexampleguided,le2025verificationguided}. MPC-CBF and sim-to-real correction methods address safety, feasibility, or transfer through control and correction layers \cite{zeng2021mpccbf,jiang2024transic}. BodyShield is distinct only in the bounded sense used here: it searches embodiment-control perturbations, repairs against discovered failures, and demands held-out and hardware-gated evidence before final claims.
\section{Problem Formulation}
Let \(z\in\mathcal{Z}\) be an embodiment-control perturbation and \(c(z)\) a normalized cost. For policy \(\pi_\theta\), task \(\tau\), and robot archetype \(r\), BodyBreak estimates \(z^\star=\arg\min c(z)\) subject to \(S(\pi_\theta,\tau,r,z)\leq\alpha\) under a finite evaluator budget. BodyShield then optimizes a repaired policy over the discovered break set, training perturbations, and near-boundary cases.
\section{BodyBreak Method}
BodyBreak compares random, one-axis, grid, and compound adversarial search. It reports evaluator-call budgets, estimated breaking cost, threshold sensitivity, and dense post-hoc challenges. It does not claim global minimality.
\section{BodyShield Repair Method}
BodyShield uses discovered failure axes to allocate repair capacity, then evaluates nominal retention, seen perturbations, held-out perturbation families, physical-style proxy shifts, secondary costs, and oracle feasibility. The current policy family is analytic and CPU-only; it validates the pipeline, not hardware deployment.
\section{Simulation Experiments}
The local run evaluates eight tasks, six robot/body archetypes, ten policy families, major embodiment-control perturbation families, compound perturbations, confidence intervals, and threshold sensitivity. Table~\ref{tab:analytic-success} summarizes the main analytic buckets. Evidence lives in \path{results/}, \path{logs/sim/results.jsonl}, \path{tables/sim_main_results.csv}, and \path{tables/sim_budget_matched_results.csv}.
\begin{table}[t]
\caption{Analytic success rates by perturbation bucket.}
\label{tab:analytic-success}
\centering\footnotesize
\begin{tabular}{lccc}
\toprule
Method & Nominal & Seen & Held-out \\
\midrule
Nominal & 0.872 & 0.757 & 0.761 \\
Domain rand. & 0.844 & 0.748 & 0.742 \\
BodyShield & 0.859 & 0.800 & 0.792 \\
Oracle & 0.925 & 0.910 & 0.907 \\
\bottomrule
\end{tabular}
\end{table}
\section{Real Robot Experiments}
\textbf{Hardware placeholder only. Do not fill until SO-ARM101/SO-101 safety gates pass.} Required hardware evidence includes noise floor, verifier agreement, reset reliability, emergency-stop test, all-trials logs, videos with failures and ambiguous cases, held-out physical modifications, and oracle feasibility.
\section{EPEC and Human-Effect Stress Test}
Human/effect and EPEC-style policies are included only as stress-test alternatives. They are not the main method and are not used to claim human-video or foundation-policy superiority.
\section{Ablations}
Ablations cover search mode, repair baseline, threshold, secondary execution costs, and conservative behavior. They are logged in \path{results/breaking_search.csv}, \path{results/threshold_sensitivity.csv}, and \path{results/secondary_metrics_by_method.csv}.
\section{Failure Analysis}
Failure categories include tracking error, calibration error, joint limit, collision, slip, unreachable pose, verifier uncertainty, and mechanical proxy faults. The paper must keep these tied to analytic simulation until hardware logs exist.
\section{Limitations}
The current package lacks hardware noise floor, camera verifier accuracy, reset reliability, real held-out physical modifications, all-trials hardware videos, external trained-policy rollouts, and real corrective-trace adaptation.
\section{Reproducibility Statement}
The package includes tests, CI, manifests, release ZIP, claim ledger, citation verification, and `python -m bodyshield.analysis.verify_package --json`. Trial fields follow \path{data_schema.json}.
\section{Ethics and Safety Statement}
The repository forbids raw hardware command paths and exposes only bounded primitives that refuse to run before explicit physical safety confirmation. This does not guarantee hardware safety.
\bibliographystyle{IEEEtran}
\bibliography{references}
\end{document}
